\section{Limitations}\label{sec:limits}

\paragraph{X25519 is a prime-order group in the models.} Clamping, the
cofactor, the twist, low-order points and non-canonical encodings are not
modelled. RFC~7748 decoding is many-to-one: the top bit is ignored and
values at least $p$ are reduced, so several byte strings name the same
point. X-Wing hashes the raw bytes of $\ct_X$, so distinct encodings give
distinct keys, but the models, whose ciphertexts are group elements, cannot
express a re-encoded ciphertext at all; Kang, Lee and Son study why such
public values must be hashed, with X-Wing as their case
study~\cite{KLS26}. \Cref{sec:lowcaught} shows that this gap is where the
implementation diverged. A cofactor-aware model, as in CryptoVerif's
WireGuard analysis~\cite{LBB19}, which uses the library's X25519-specific
Diffie--Hellman macros, would close part of it. We have not attempted one.

\paragraph{Composition is by hand.} CryptoVerif checks each model
separately. Substituting \cref{lem:kem,lem:sig} into
\cref{thm:auth,thm:secrecy}, the KEM side conditions, and the argument that
SHA3-256's three idealisations do not interfere are all done on paper
(\cref{sec:status}). So is the step from per-role multi-test secrecy to the
single-Test game of the original sketch.

\paragraph{Remaining assumptions.} The ML-KEM leg of the ML-KEM-1024 profile
still rests on a split-key PRF assumption for HKDF, which no published result
covers. The protocol's KDF rests on HMAC's swap-PRF property for a key
distribution whose feasibility in the sense of~\cite{BBGS23} we have argued
but not checked.

\paragraph{The adversary.} Ephemeral-key reveal is static. Session-state
reveal and reveal of the responder's encapsulation coins are not modelled. A
fresh session's key can be tested but not revealed in place of a test. There
is one ciphersuite, so cross-suite downgrade resistance rests on the
symbolic model alone.

\paragraph{The bounds are not tight.} Reduced to single-instance IND-CCA2,
the KEM term loses a factor of about $n_S^2$ (\cref{sec:reviewcaught}). We
have not computed concrete security levels.

\paragraph{Models, specification and code.} The models follow the
specification, and their correspondence with it and with the Python
implementation was checked by inspection. The implementation itself is not
verified. It is research code that trusts pyca/cryptography for every
primitive, and we make no claim about side channels. The proof scripts name
variables that CryptoVerif generates and count occurrences, so they are tied
to CryptoVerif~2.13 and break when a model changes.

\paragraph{Review.} The development was AI-assisted, and its independent
review was done by a separate instance of the same AI system. That review
found real problems, but it shares blind spots with the development that a
human expert might not. Neither the models nor this paper has had external
peer review.

\paragraph{Scope.} HYBRID-AKE is an engineering design, not a standard. For
production systems, TLS~1.3 with a standardised hybrid group from a
maintained library is the better choice. The results here are about one
design and one implementation, and the lessons in \cref{sec:lessons} are
offered as observations from it, not as general findings.
